% ============================================================
%  part3/ch15-allen-cahn.tex  —  Chapter 15: Allen-Cahn Dynamics
% ============================================================

\chapter{Allen--Cahn Dynamics}
\label{ch:allen-cahn}

\epigraph{%
  Everything descends.%
}{---}

\section{From discrete to continuum}

The Ising model is discrete.
Allen--Cahn is its continuum counterpart.
Replace spins by a scalar field $\phi(x,t) \in \mathbb{R}$.

\begin{definition}[Ginzburg--Landau free energy]
  \[
    \mathcal{F}[\phi]
    = \int \left[\frac{\epsilon^2}{2}|\nabla\phi|^2 + W(\phi)\right]dx,
    \quad W(\phi) = \frac{1}{4}(\phi^2-1)^2.
  \]
\end{definition}

The gradient term penalises interfaces.
The potential $W$ favours $\phi = \pm 1$ (the two phases).

\begin{proposition}[Euler--Lagrange equation]
  Critical points of $\mathcal{F}$ satisfy
  \[
    -\epsilon^2\Delta\phi + W'(\phi) = 0.
  \]
\end{proposition}

\begin{proof}
  Take variation $\phi \mapsto \phi + \varepsilon\psi$:
  \[
    \frac{d}{d\varepsilon}\mathcal{F}
    = \int\left[\epsilon^2\nabla\phi\cdot\nabla\psi
      + W'(\phi)\psi\right]dx
    = \int\left[-\epsilon^2\Delta\phi + W'(\phi)\right]\psi\,dx.
  \]
  Since $\psi$ is arbitrary, the bracket vanishes.
\end{proof}

\section{The Allen--Cahn equation}

\begin{definition}[Allen--Cahn equation]
  \[
    \boxed{\partial_t\phi = M(\epsilon^2\Delta\phi - W'(\phi)).}
  \]
\end{definition}

\begin{theorem}[Energy dissipation]
  \label{thm:ac-dissipation}
  \statusproven{}
  Solutions of Allen--Cahn satisfy
  \[
    \frac{d\mathcal{F}}{dt}
    = -M\int\left(\frac{\delta\mathcal{F}}{\delta\phi}\right)^2 dx
    \leq 0.
  \]
\end{theorem}

\begin{proof}
  $\frac{d\mathcal{F}}{dt}
  = \int\frac{\delta\mathcal{F}}{\delta\phi}\partial_t\phi\,dx
  = -M\int\left(\frac{\delta\mathcal{F}}{\delta\phi}\right)^2 dx
  \leq 0$.
\end{proof}

\begin{StatusBox}{Proven Framework}
  Theorem~\ref{thm:ac-dissipation} is the template for all
  free-energy dissipation results in this book.
  The RSVP dissipation theorem in
  Chapter~\ref{ch:energy-dissipation} is its direct descendant.
\end{StatusBox}

\section{Interface limit and coarsening}

In the limit $\epsilon \to 0$, interfaces become thin.
Matched asymptotic analysis shows Allen--Cahn reduces to
curvature flow: $V_n = -M\kappa$.

\begin{corollary}
  Allen--Cahn domains obey $\ell(t) \sim t^{1/2}$, $z = 2$.
\end{corollary}

\section{Dynamic scaling}

\begin{definition}[Two-point correlation function]
  \[
    C(r,t) = \langle\phi(x,t)\phi(x+r,t)\rangle.
  \]
\end{definition}

\begin{definition}[Dynamic scaling hypothesis]
  \[
    C(r,t) = f\!\left(\frac{r}{\ell(t)}\right).
  \]
\end{definition}

Under dynamic scaling the coarsening system is characterised
by a single growing length scale $\ell(t)$.
Data from different times collapse onto one curve when plotted
against $r/\ell(t)$.

\begin{ForwardRef}
  The correlation function $C(r,t)$ defined here will become
  the primary observable for testing RSVP predictions in
  Chapter~\ref{ch:correlation-functions}.
  The coarsening exponent $z = 2$ will be compared to
  $z_{\mathrm{RSVP}} = z(Pe, \Gamma, \Omega)$ conjectured in
  Chapter~\ref{ch:open-problems}.
\end{ForwardRef}

\begin{CrossDomain}[Allen--Cahn as template for RSVP]
  \begin{center}
  \begin{tabular}{@{}ll@{}}
    \toprule
    \textbf{Allen--Cahn} & \textbf{RSVP} \\
    \midrule
    Scalar field $\phi$ & Plenum field $\Phi$ \\
    Double-well $W(\phi)$ & Lamphron potential $U(\ell)$ \\
    Interface energy $|\nabla\phi|^2$ & Admissibility energy \\
    Curvature flow & Entropy descent \\
    Domain walls & Filaments and void boundaries \\
    $d\mathcal{F}/dt \leq 0$ & $d\mathcal{F}_{\mathrm{RSVP}}/dt \leq 0$ \\
    \bottomrule
  \end{tabular}
  \end{center}
\end{CrossDomain}

\WhatSurvived{Allen--Cahn coarsening}{%
  Phase assignment ($\pm 1$)\\
  Domain topology\\
  Energy (decreasing)%
}{%
  Interface positions\\
  Sub-domain fluctuations\\
  Initial field values%
}{%
  Phase identity recoverable.\\
  Interface history lost.%
}{%
  $\ker$: all initial conditions\\
  with the same asymptotic\\
  phase assignment.%
}
